\documentclass[11pt]{article}

\usepackage{maiacustom}

\begin{document}

\psettitle{Astronomy Question Bank}

These questions were carefully written or selected to prepare students for the astronomy team selection process in Brazil. Some questions are not original; those are marked with () before the statement. The question-bank template is the same one used by Professor Kevin Zhou. His work is invaluable, and many ideas in this set can be found in his handouts.

% \begin{psidea}{Idea title}{}
% Idea
% \end{psidea}

% \begin{psexample}{Example title}{}
% Example
% \end{psexample}

% \begin{pssolution*}{}{}
% Solution
% \end{pssolution*}

% \begin{psremark*}{Remark title}{}
% Remark
% \end{psremark*}
\section{Optics and Telescopes}

\pts{3}
    \begin{pproblem} A Keplerian telescope has two converging lenses. The first has a radius of curvature of both faces equal to \(R = 2\text{ m}\), and the second has both radii of curvature equal to \(r = 0.5 \text{ m}\). Considering that both lenses have negligible thickness and are made of a material with refractive index \(n=1.4\), calculate the diameter and the magnification of the telescope, knowing that this telescope is an f/12.
    \begin{pssolution*}{}{}
        The focal length of the telescope is the sum of the individual focal lengths of each lens. Using the lensmaker's equation,
        \[\frac{1}{f} = \left(\frac{n_{\text{lens}}}{n_{\text{medium}}}-1\right)\left(\frac{1}{R_1}+\frac{1}{R_2}\right)\]

        Let \(a\) and \(b\) be the objective and eyepiece lenses, respectively. We obtain
        
        \[f_a = 2.50 \text{ m}\ ; \ f_b = 0.625\text{ m}\]

        Therefore, the focal length of the telescope is \[f = f_a+f_b = 3.125\text{ m}\]. Using that it is an \(f/12\), we obtain
        \[\boxed{D = f/12 = 260 \text{ mm}}\]
    \end{pssolution*}
    \end{pproblem}


\pts{5}
    \begin{pproblem}
        In this problem, we will become familiar with one of the most powerful tools in optics, \textit{geometric optics}. The goal of this tool is to model lenses and mirrors in the form of matrices. This is very useful for solving problems involving combinations of several lenses, and it is a method for solving geometric-optics problems (practically) without using geometry. For that, pay attention to the following definitions in the figure below.

        \begin{figure}[H]
            \centering
            \includegraphics[width=0.9\linewidth]{imagens/lentes1.png}
            \caption{Diagram 1}
        \end{figure}

        All dashed lines are parallel to the optical axis, represented by the horizontal arrow. \(P_k\) denotes the point where the light passes from one medium to another, and \(y_k\) is the vertical coordinate of the point \(P_k\). The subscripts \(i\) and \(t\) are used to refer to incident and transmitted, respectively. Thus, for example, \(\alpha_{i,k}\) is the angle that the incident light ray makes with the horizontal at the point \(P_k\), while the angle \(\alpha_k\) is the angle between the point \(P_k\) and the center of lens \(k\). For example, in the figure, the angle \(\alpha_1\) is represented by \(V_1\hat{C}P_1\).
        

        For this problem, we will assume that all angles of interest are small, so that \(\tan\theta \approx \sin\theta \approx \theta\).
        \begin{alternativas}
            \item Starting from Snell's law, find a relation between \(n_{i,1}\), \(n_{t,1}\), \(\alpha_1\), \(\alpha_{i,1}\), and \(\alpha_{t,1}\).
        
            \item Find a relation for \((1) \ n_{t,1}\alpha_{t,1}\) and \((2) \ y_{t,1}\). Leave your answers in terms of \(n_{i,1}\), \(\alpha_{i,1}\), \(y_{i,1}\), and \(\mathcal{D}_1\), where
            \[\mathcal{D}_1 = \frac{n_{t,1}-n_{i,1}}{R_1}\]

        Now for one more definition. Let the vector \(\mathbf{r}_{t,1}= (n_{t,1}\alpha_{t,1}, \ y_{t,1})\) and \(\mathbf{r}_{i,1}= (n_{i,1}\alpha_{i,1}, \ y_{i,1})\).
        
            \item The two equations found in the previous item can be written in matrix form, as
            \[\mathbf{r}_{t,1}=\mathcal{R}_1\mathbf{r}_{i,1}\]
            Find the \(2\times 2\) matrix equivalent to \(\mathcal{R}_1\).

        Our interest now is to find the relations between the points \(P_2\) and \(P_1\).

            \item Find an equation for \(n_1\alpha_{i,2}\) and \(y_{i,2}\). Using the reasoning of the previous item, leave your answer in the form
            \[\mathbf{r}_{i,2}=\Gamma_{2,1}\mathbf{r}_{t,1}\]
            where \(\Gamma_{2,1}\) is also a \(2\times 2\) matrix. Leave your answer in terms of the lens thickness, \(d_{2,1}\).

            \item We define the lens matrix, \(\mathcal{A}_{2,1}\), by the following equation:
            \[\mathbf{r}_{t,2} = \mathcal{A}_{2,1}\mathbf{r}_{i,1}\]
            so that

            \[\mathcal{A}_{2,1} = \begin{pmatrix}
                a_{11} & a_{12} \\
                a_{21} & a_{22} \\
            \end{pmatrix}\]

            Find \(\mathcal{A}_{2,1}\) explicitly. This is extremely important, because it describes the light ray that leaves the lens in terms of the ray that enters.

            \item Among all the properties of the lens matrix, the most curious is that the term \(a_{12}\) is proportional to \(-1/f\), where \(f\) is the focal length of the lens. Prove this result. (Hint: you will be able to write \(-a_{12}\) with a well-known equation from optics.)
        
        Now let us get to work and make practical use of matrix optics.
        
            \item Consider the following diagram:
                \begin{figure}[H]
                    \centering
                    \includegraphics[width=0.9\linewidth]{imagens/lentes2.png}
                    \caption{Diagram 2}
                \end{figure}
                 Find \(y_i\) in terms of \(y_0\) and the data in the figure.
                 
            \item If the lens is between two different media, \(n_{i,1} \ne n_{t,2}\), we have the following relation:
            \[a_{12} = -\frac{n_{i,1}}{f_0}=-\frac{n_{t,2}}{f_I}\]
                Using this, redo the previous item, considering that there is water with \(n_w = 4/3\) between the two lenses.
        
            \item A very common lens arrangement is the Tessar objective, found in many cameras because of its efficiency in reducing aberration and astigmatism. The arrangement has the following form:
                \begin{figure}[H]
                    \centering
                    \includegraphics[width=0.7\linewidth]{imagens/lentes3.png}
                    \caption{Tessar objective diagram}
                \end{figure}
                Find the equivalent lens matrix, \(\mathcal{A}_{1,7}\), in terms of \(\mathcal{R}_i\) and \(\Gamma_{i,j}\).
        \end{alternativas}
    \begin{pssolution*}{}{}
        \begin{alternativas}
            \item Snell's law tells us that \(n_{i,1}\theta_{i,1} = n_{t,1}\theta_{t,1}\), where \(\theta\) are the angles that the light makes with the normal. Note that we can write both angles \(\theta\) as
            \[\theta_{i,1} = \alpha_1 + \alpha_{i,1}, \ \ \theta _{t,1} = \alpha_1 + \alpha_{t,1}\]

            Therefore, the relation we are looking for is

            \[\boxed{n_{i,1}(\alpha_1+\alpha_{i,1}) = n_{t,1}(\alpha_1+\alpha_{t,1})}\]

            \item Substituting \(\alpha_1 = y_1/R_1\), we have
            
            \[n_{i,1}(y_1/R_1+\alpha_{i,1})=n_{t,1}(y_1/R_1 +\alpha_{t,1})\]

            Isolating \(n_{t,1}\alpha_{t,1}\),

            \[n_{t,1}\alpha_{t,1} = n_{i,1}\alpha_{i,1}+ \frac{y_1}{R_1}(n_{i,1}-n_{t,1})\]

            In terms of \(\mathcal{D}_1\),

            \[n_{t,1}\alpha_{t,1} = n_{i,1}\alpha_{i,1} - y_1\mathcal{D}_1\]

            Note that \(y\) does not change immediately after the light is transmitted. Therefore \(y_1\equiv y_{i,1}\equiv y_{t,1}\).

            \item In matrix form,
            
            \[\left(\begin{matrix}
                n_{t,1}\alpha_{t,1} \\
                y_{t,1}\\
            \end{matrix}\right) = \left(\begin{matrix}
                1 & - \mathcal{D}_1 \\
                0 & 1
            \end{matrix}\right)\left(\begin{matrix}
                n_{i,1}\alpha_{i,1} \\
                y_{i,1}\\
            \end{matrix}\right)\]

            Thus, our matrix \(\mathcal{R}_1\), also known as the refraction matrix, is given by

            \[\boxed{\mathcal{R}_1 = \left(\begin{matrix}
            1 & -\mathcal{D}_1 \\
            0 & 1 \\
            \end{matrix}\right)}\]

        \item First, let us find the height of point \(2\). Let \(d_{21}\) be the thickness of the lens. Using trigonometry, we obtain
        
        \[y_{i,2} = y_{t,1} + \alpha_{t,1}d_{21}\]

        We have that \(\alpha_{t,1}\) and \(\alpha_{i,2}\) are alternate interior angles. That is, \(\alpha_{t,1}=\alpha_{i,2}\). Since the refractive index inside the lens is constant (\(n_{i,2} = n_{t,1}\)), we have \(n_{t,1}\alpha_{t,1} = n_{i,2}\alpha_{i,2}\). Putting the two expressions in matrix form, we obtain

        \[\begin{pmatrix}
            n_{i,2}\alpha_{i,2}\\
            y_{i,2}\\
        \end{pmatrix}
        =
        \begin{pmatrix}
            1 & 0 \\
            \frac{d_{2,1}}{n_{t,1}} & 1 \\
        \end{pmatrix}
        \begin{pmatrix}
            n_{t,1}\alpha_{t,1} \\
            y_{t,1}\\
        \end{pmatrix}
        \]

        Therefore, we can conclude that

        \[\boxed{\Gamma_{2,1} = \begin{pmatrix}
        1 & 0 \\
        \frac{d_{2,1}}{n_{t,1}} & 1 \\
        \end{pmatrix}}\]

        \item Combining the expressions found earlier,
        \[\mathbf{r}_{i,2} = \Gamma_{2,1} \mathbf{r}_{t,1}\]
        \[\mathbf{r}_{i,2} = \Gamma_{2,1} (\mathcal{R}_1\mathbf{r}_{i,1})\]

        And, as we saw earlier, to go from \(\mathbf{r}_{i,k}\) to \(\mathbf{r}_{t,k}\) it is enough to use the matrix \(\mathcal{R}_k\). Thus,

        \[\mathbf{r}_{t,2} = \mathcal{R}_2\Gamma_{2,1}\mathcal{R}_1\mathbf{r}_{i,2}\]

        Therefore, the matrix \(\mathcal{A}_{2,1}\) is the product of the matrix \(\mathcal{R}_2\), the matrix \(\Gamma_{2,1}\), and the matrix \(\mathcal{R}_1\). Explicitly,

        \[\mathcal{A}_{2,1} = \left(\begin{matrix}
            1 & -\mathcal{D}_2 \\
            0 & 1 \\
            \end{matrix}\right)\begin{pmatrix}
        1 & 0 \\
        \frac{d_{2,1}}{n_{t,1}} & 1 \\
        \end{pmatrix}
        \left(\begin{matrix}
            1 & -\mathcal{D}_1 \\
            0 & 1 \\
            \end{matrix}\right)\]
        \end{alternativas}

        With enough determination, we can carry out this calculation, obtaining

        \[\boxed{
        \mathcal{A}_{2,1} = 
        \begin{pmatrix}
        1 - \frac{\mathcal{D}_2d_{2,1}}{n_{t,1}} & - \mathcal{D}_1 - \mathcal{D}_2 + \frac{\mathcal{D}_1\mathcal{D}_2d_{2,1}}{n_{t,1}} \\
        \frac{d_{2,1}}{n_{t,1}} & \frac{\mathcal{D}_1 d_{2,1}}{n_{t,1}}\\
        \end{pmatrix}
        }
        \]

        \item Working with the term \(a_{1,2}\), we have
        
        \[a_{1,2} = -\mathcal{D}_1 - \mathcal{D}_2 + \frac{\mathcal{D}_1\mathcal{D}_2d_{2,1}}{n_{t,1}}\]

        Substituting \(\mathcal{D}_k\), we have

        \[a_{1,2} = -\frac{n_{t,1}-n_{i,1}}{R_1} - \frac{n_{t,2}-n_{i,2}}{R_2} + \frac{(n_{t,1}-n_{i,1})(n_{t,2}-n_{i,2})d_{2,1}}{n_{t,1}R_1R_2}\]

        Note that, in our situation, \(n_{i,1}\) and \(n_{t,2}\) correspond to the refractive indices of the medium, while \(n_{t,1}\) and \(n_{i,2}\) correspond to the refractive indices of the lens, which will be denoted by \(n_m\) and \(n_l\), respectively. Thus,

        \[a_{1,2} = -\frac{n_{l}-n_{m}}{R_1} - \frac{n_{m}-n_{l}}{R_2} + \frac{(n_{l}-n_{m})(n_{m}-n_{l})d_{2,1}}{n_lR_1R_2}\]

        \[a_{1,2} = - (n_l-n_m)\left(\frac{1}{R_1}-\frac{1}{R_2}\right) - \frac{(n_l-n_m)^2 d_{2,1}}{n_lR_1R_2}\]

        \[a_{1,2} = -(n_l-n_m)\left(\frac{1}{R_1}-\frac{1}{R_2} + \frac{(n_l-n_m)d_{2,1}}{n_l R_1R_2}\right)\]

        This has the same form as the lensmaker's equation (for non-negligible thicknesses).

        \[\frac{1}{f} = \left(\frac{n_l}{n_m}-1\right)\left(\frac{1}{R_1}-\frac{1}{R_2} + \frac{(n_l-n_m)d_{2,1}}{n_l R_1R_2}\right)\]

        Thus, we can conclude that

        \[\boxed{a_{1,2} = -\frac{n_m}{f}}\]

\end{pssolution*}

\end{pproblem}

\pts{2}
\begin{pproblem}
    The wave theory of light shows that a perfect focus is not possible, because of diffraction effects associated with the finite aperture of the lens. This lack of a perfect focus prevents objects that are very close together from being distinguished. This problem can be studied from two different points of view:

    The wave theory of light predicts that a lens of diameter $D$ cannot focus a parallel beam of light of wavelength $\lambda$ into an angle smaller than the diffraction limit:
    \begin{equation}
        \theta_m \approx 1.22 \frac{\lambda}{D}
    \end{equation}

    Now consider a quantum approach: the photons that are focused by the lens. These photons are known to have passed somewhere within a radius of the center of the lens. The uncertainty in the position $x$ is associated with an uncertainty in the $x$ component of the photon's momentum. Consequently, a photon that, in the absence of this uncertainty, would have been brought onto the optical axis of the focal plane may now be deflected by an angle $\theta \ll 1$.

    Consider the de Broglie wavelength $\lambda = \frac{h}{p}$. Find a limit for $\theta$.

    \textbf{Hint:} Heisenberg's uncertainty principle relates the imprecision between measurements of the momentum and the position of a particle by
    \[\Delta p_x \Delta x = \frac{\hbar}{2}\]

\begin{pssolution*}{}{}
    The momentum of a photon is given by \(p = h/\lambda\), and we know that the photon has \(\Delta x = D\). Using Heisenberg's uncertainty principle,

    \[\Delta p_x = \frac{\hbar}{2D}\]

    Using \(\theta \approx \frac{\Delta p_x}{p}\), we obtain
    \[\boxed{\theta \approx\frac{\lambda\hbar}{2D h} = \frac{1}{4\pi}\frac{\lambda}{D}}\]
\end{pssolution*}
\end{pproblem}

\pts{3}
\begin{pproblem} (Magna notes)
    Prove that the quantity \(n_iR_i\sin \theta_i\), for spherical shells adjacent to one another, is an invariant, where \(n_i\) is the refractive index of the \(i\)-th spherical shell, \(\theta_i\) is the angle that the light ray makes with the normal to the \(i\)-th spherical shell, and \(R_i\) is the radius of the \(i\)-th spherical shell.
    
\end{pproblem}

\pts{4}
\begin{pproblem}
    The refractive index of a planet's atmosphere is given by

    \[n(h) = \frac{n_0}{1+\epsilon h}\]

    where \(n_0\) and \(\epsilon\) are constants.

    \begin{alternativas}
        \item A light ray reaches the atmosphere parallel to the surface, at a height \(h'\ll R\). What is the trajectory?
        \item Given the trajectory of the light ray, calculate the radius of the planet.
    
        \textbf{Hint:} You may find the following relation useful,

        \[\int \frac{t}{\sqrt{a-bt^2}}dt = -\frac{1}{b}\sqrt{a-bt^2}+C\]
    \end{alternativas}
\begin{pssolution*}{}{}
    \begin{alternativas}
        \item Defining \(ds = \sqrt{dx^2+dh^2}\) as the infinitesimal arc-length element, we have
        \[\sin\theta = \frac{dx}{ds}, \ \ \cos\theta = \frac{dh}{ds}\]

        Since the medium is continuous, we can write

        \[n(h)\sin\theta(h)=\mathrm{const}\]

        Thus,

        \[n\frac{dx}{ds} = K\]

        Substituting \(n\),

        \[\frac{dx}{ds} = \frac{K(1+\epsilon h)}{n_0}\]

        However, by the definition of \(s\), we have

        \[\left(\frac{dx}{ds}\right)^2 + \left(\frac{dh}{ds}\right)^2 = 1\]

        Substituting \(dx/ds\),

        \[\frac{dh}{ds} = \sqrt{1- \left(\frac{K(1+\epsilon h)}{n_0}\right)^2}\]

        Using the chain rule,

        \[\frac{dh}{dx} = \frac{dh/ds}{dx/ds} = \frac{\sqrt{1- \left(\frac{K(1+\epsilon h)}{n_0}\right)^2}}{\frac{K(1+\epsilon h)}{n_0}}\]

        Simplifying the expression,

        \[\frac{dh}{dx} = \frac{\sqrt{n_0^2 - K^2(1+\epsilon h)^2}}{K(1+\epsilon h)}\]

        To solve this ODE, let us introduce the variable \(u = 1+\epsilon h\), so that

        \[dh = \frac{du}{\epsilon}\]

        Thus,

        \[\frac{du}{dx} = \epsilon\frac{\sqrt{n_0^2 - K^2u^2}}{Ku}\]

        Separating the terms and integrating,

        \[\int \frac{u}{\sqrt{n_0^2 - K^2u^2}} du = \frac{\epsilon}{K}\int dx\]
    
        The integral on the left-hand side is the same one given in the problem, so

        \[-\frac{1}{K^2}\sqrt{n_0^2-K^2u^2} = \frac{\epsilon}{K}x + C\]

        Working with this expression,

        \[\sqrt{n_0^2-K^2u^2} = -\epsilon K x - K^2C\]

        \[n_0^2-K^2u^2 = (\epsilon K x + K^2C)^2\]

        Substituting \(u\),

        \[n_0^2-K^2(1+\epsilon h)^2 = \epsilon^2K^2x^2 + 2\epsilon K^3 C x + K^4C^2\]

        Rearranging the terms,

        \[\epsilon^2K^2x^2 + 2\epsilon K^3 C x + K^2(1+\epsilon h)^2 = n_0^2 - K^4C^2 \]

        \[\epsilon^2x^2 + 2\epsilon K C x + (1+\epsilon h)^2 = n_0^2/K^2 - K^2C^2 \]

        Dividing by \(\epsilon^2\),

        \[x^2 + \frac{2KC}{\epsilon}x + \frac{K^2C^2}{\epsilon^2} + \left(h+\frac{1}{\epsilon}\right)^2 = \left(\frac{n_0}{\epsilon K}\right)^2\]

        Collecting the perfect square,

        \[\boxed{\left(x + \frac{KC}{\epsilon}\right)^2+\left(h+\frac{1}{\epsilon}\right)^2 = \left(\frac{n_0}{\epsilon K}\right)^2}\]
   
        which is precisely the equation of a circle of radius \(\frac{n_0}{\epsilon K}\) and center \(C = (-\frac{KC}{\epsilon}, \ - \frac{1}{\epsilon})\).

        \item In the case of a ray that is tangent to the planet, that ray has \(\theta_0 = \pi/2\) and \(h_0=0\). Thus, using
        
        \[n(h)\sin\theta(h) = n_0\sin\theta_0 = K\]

        we obtain \(K = n_0\). Thus, the radius of the planet is given by the radius of the trajectory (the light ray is tangent to the surface) and, therefore,

        \[\boxed{R = \frac{1}{\epsilon}}\]
    \end{alternativas}
    
\end{pssolution*}
\end{pproblem}

\end{document}
